\documentclass[book.tex]{subfiles}

\begin{document}

\chapter{Autoencoders and Variational Autoencoders}
\label{chap:autoencoders}
\noindent\rule{11cm}{0.4pt} \\
\textbf{Prerequisites:}  \autoref{chap:ml_basics}, \autoref{chap:grad_descent}  \\
\textbf{Difficulty Level:} ** \\
\noindent\rule{11cm}{0.4pt}
\vspace{5mm}

Often in scientific machine learning applications, we lack any labeled training data. We may have a collection of molecular structures or materials formulas, but may not have access to any associated experimental labels for these data points. Can we still do machine learning?

\section{Autoencoder Reconstruction}

Suppose the input $x$ is in space $\mathcal{S}$. An encoder is a function $f_{\theta}: \mathcal{S} \to \mathbb{R}^n$ that transforms elements of $\mathcal{S}$ into $n$ dimensional vectors parameterized by $\theta$. Let $g_{\theta}: \mathbb{R}^n \to \mathcal{S}$ be a map in the opposite direction. An autoencoder posits that the encoding and decoding an input can be a meaningful tool to force learning of the structure of a dataset; the challenge of representation learning is picking $\theta$ such that 
\begin{align}
    \theta^* = \argmin_{\theta} \|g_{\theta}(f_{\theta}(s)) - s \| 
\end{align}

One way to think about $g_{\theta} \circ f_{\theta}$ is as a learned version of the identity. There are a couple of reasons learning the identity may be useful. Suppose that 
\begin{align}\mathcal{S} = \mathbb{R}^M,\quad M \gg n
\end{align} 
In this case, $g_{\theta} \circ f_{\theta}$ is a compressed version of the identity that mandates an internal compression to $n$-dimensional space.
\begin{figure}
    \centering
    % \includegraphics{figures/Differentiable Models/autoencoders/autoencoder_compression.png}
    \includegraphics{figures/Differentiable Models/autoencoders/Autoencoder.png}
    %\caption{\textcolor{red}{TODO: Replace. From \url{https://blog.paperspace.com/autoencoder-image-compression-keras/}}}
    \caption{Autoencoders encode the input into a suitable embedding space and decode the embedded vector into the output space. This encoding and decoding process can help networks learn about the meaningful structure of the data.}
    \label{fig:auto_compression}
\end{figure}
As a slightly more complex example, we may have $\mathcal{S} = \mathbb{R}^{a \times b \times c}$ be a space of images with $a \times b$ pixels and $c$ colors. In this case, the encoder $f_{\theta}$ transforms these images into $n$-dimensional feature vectors.

We can choose different architectures to encode and decode, but a common simple choice is to just use fully connected networks.

\begin{lstlisting}[language=python]
class Autoencoder(f: FullyConnectedNetwork, g: FullyConnectedNetwork):
  def $\lambda$(X: $\mathbb{R}$[N, T, d])) $\to$ $\mathbb{R}$[N, T, d]): 
    g(f(X))
    
def reconstruction_loss(model, X):
  return $\|$model(X) - X$\|^2$
\end{lstlisting}

\section{Variational Autoencoder}

A variational autencoder can be viewed as a classical autoencoder with the addition of a regularization term to the loss. Here we train the encoder $f$ to produce the means and standard deviations of the encoded distribution.
\begin{align}
\mu_x, \sigma_x &= f(x)
\end{align}

Then during decoding, we sample from this distribution
\begin{align}
    z \sim \mathcal{N}(\mu_x, \sigma_x)
\end{align}
And then reconstruct the output.
\begin{align}
    \mathrm{out} &= g(z)
\end{align}

Converting into Physika, we obtain

\begin{lstlisting}[language=python]
class VariationalAutoencoder(f: FullyConnectedNetwork, g: FullyConnectedNetwork):
  def $\lambda$(X):
    $\mu_x$, $\sigma_x$ = f(X)
    z = sample($\mathcal{N}$($\mu_x$, $\sigma_x$))
    g(z)

\end{lstlisting}

Note that the variational autoencoder provides a natural mechanism to draw samples from a learned distribution by sampling random $z$ and then decoding accoding to $g$. The latent distribution is controlled by adding a KL-divergence. %(\textcolor{red}{TODO: Add KL-divergenece to the preliminaries}).

\begin{align}
    \mathcal{L} &= \|x - f(\textrm{sample}(g(x)))\|^2 + \mathrm{KL}(\mathcal{N}(\mu_x, \sigma_x), \mathcal{N}(0, I))
\end{align}

The gradient descent step here is more complex since we have to do a backwards update over a sampling step. Empirically, we can compute this quantity in expectation, but proofs of correctness are quite tricky for these probabilistic losses.
%\url{https://towardsdatascience.com/understanding-variational-autoencoders-vaes-f70510919f73}
\end{document}